% ---------------------------------------------------------------------
%  7. Discussion: when cooperation helps + improvements
% ---------------------------------------------------------------------
\section{Discussion: when does cooperation help?}
\label{sec:discussion}

The study gives a specific answer, and part of it is a clean ``not here.''

\textbf{Routing over complementary strengths helps.} The two models fail on
different problems (Finding~1), so a portfolio that tries the free tactics, then
one model's repair, then the other's, captures the union of their abilities in one
policy --- which is why the integrated agent (5/9) clears the vanilla solos (3/9)
and their union (4/9).

\textbf{Same-model repair helps.} Feeding a model its own failed proof plus the
exact Lean diagnostic is the lever that unlocks \texttt{p10} and supplies the
margin above the union (Finding~2). Notably, this needs only one model at a time.

\textbf{Cross-model handoff does not help, and can hurt.} This is the negative
result of Finding~3: handing one model's proof to the other for repair added no
win and, on \texttt{p05}, destroyed one. The honest conclusion is that the value in
this system comes from \emph{verification and routing}, not from the two models
repairing each other. I keep handoff in the ladder only as a late, seeded,
low-cost step, because on a different problem distribution it might occasionally
help --- but on this evidence it is a residual tool, not the point.

\textbf{On the hard tail, capability comes from search, not conversation.} The
\texttt{p09} win is test-time sampling scale (\S\ref{sec:frontier}); the
hollow-bank boundary shows coordination cannot conjure an assembled proof from
verified fragments when neither model can direct the assembly. Cooperation
reallocates and repairs capability; scale extends it; neither invents it.

\subsection{Improvements the results point to}
Each failure names a next step, all within the two-model, no-retrieval constraint:
\begin{itemize}
  \item \textbf{Measure the router baseline.} The most important open comparison is
        a same-budget policy that only routes between the two same-model repair
        arms ($\RQ\!\lor\!\RG$). If it ties the full ladder, the honest story is
        ``routing + free tactics,'' and the ladder should shed anything that does
        not earn its wall.
  \item \textbf{Budget the hard wins on purpose.} \texttt{p09} is solvable when
        enough sampling is aimed at it; a scheduler that detects a multi-theorem
        statement and front-loads per-theorem sampling would convert a
        probabilistic win into a reliable one.
  \item \textbf{Goal-shaped lemma selection and forced assembly.} The hollow bank
        fails at assembly, not generation; selecting lemmas by how they unify with
        the remaining goal, and ending each round with a mandatory full-bank close,
        attacks the actual bottleneck.
  \item \textbf{An offline hint catalogue.} Runtime retrieval is disallowed, but a
        small, generic, problem-agnostic tactic/lemma catalogue baked in before the
        network lock would give the fallback stronger closing moves without
        breaking the rules.
\end{itemize}
None of these needs a larger model or a new modality; each attacks a failure this
study exhibited.
